\documentclass{article}
\usepackage{amsmath}   % For matrices, align environments, \operatorname
\usepackage{amssymb}   % For math symbols like \oplus
\usepackage[utf8]{inputenc}  % UTF-8 support
\usepackage[T1]{fontenc}     % Font encoding
\usepackage{geometry}        % Optional: adjust margins
\geometry{margin=1in}

\title{Linear Projector: Kernel and Image}
\author{}
\date{}

\begin{document}
\maketitle

\section{Definition of a Linear Projector}
A linear projector $P$ is an idempotent operator satisfying:
\[
P^2 = P
\]

\section{\(\ker P\) (Kernel of \(P\))}
The kernel (null space) of $P$ is the set of all vectors mapped to the zero vector:
\[
\ker P = \{ \mathbf{v} \mid P\mathbf{v} = \mathbf{0} \}
\]
For a projector, this is the subspace of vectors "projected away" to nothing.

\section{\(\operatorname{im} P\) (Image of \(P\))}
The image (range) of $P$ is the set of all achievable outputs of $P$:
\[
\operatorname{im} P = \{ \mathbf{w} \mid \exists \mathbf{v}, \mathbf{w} = P\mathbf{v} \}
\]
This is the subspace that $P$ projects *onto* — all valid results of the projection live here.

\section{Key Decomposition}
The full vector space $V$ splits into a \textbf{direct sum} of these two subspaces:
\[
V = \ker P \oplus \operatorname{im} P
\]
Every vector $\mathbf{v}$ has a unique decomposition:
\[
\mathbf{v} = P\mathbf{v} + (\mathbf{v} - P\mathbf{v})
\]
where $P\mathbf{v} \in \operatorname{im} P$ and $\mathbf{v} - P\mathbf{v} \in \ker P$.

\section{Example}
Take the 2D projector:
\[
P = \begin{bmatrix}1 & 0 \\ 0 & 0\end{bmatrix}
\]
- $\operatorname{im} P = \operatorname{span}\left\{ \begin{bmatrix}1 \\ 0\end{bmatrix} \right\}$ (the x-axis)
- $\ker P = \operatorname{span}\left\{ \begin{bmatrix}0 \\ 1\end{bmatrix} \right\}$ (the y-axis)

\section{Isometries vs Unitary Operators}
\subsection{Definitions}
\begin{itemize}
    \item \textbf{Isometry} $U$: Linear operator with $U^*U = I$ (norm-preserving, injective).
    \item \textbf{Unitary operator} $U$: Linear operator with $U^*U = UU^* = I$ (norm-preserving, bijective).
\end{itemize}

\subsection{Hierarchy}
All unitary operators are isometries, but not all isometries are unitary (this distinction only appears in infinite-dimensional spaces). In finite dimensions, the two classes are identical.

\subsection{Example: Right Shift on $\ell^2$}
The right shift operator $R$:
\[
R(x_1, x_2, x_3, \dots) = (0, x_1, x_2, x_3, \dots)
\]
is an isometry ($R^*R = I$) but not unitary ($RR^* \neq I$, as it is not surjective).

\end{document}